% 原创教学示意；z 与 h 是同一隐藏层激活前后的坐标，并非两个隐藏层。
\begin{tikzpicture}[x=1mm,y=1mm,font=\figurefont\footnotesize,text=NNInk,
  bp neuron/.style={circle,minimum size=7mm,inner sep=0pt,line width=1.1pt},
  bp input/.style={bp neuron,draw=NNInput,fill=NNInput!18},
  bp hidden/.style={bp neuron,draw=NNHidden,fill=NNHidden!18},
  bp output/.style={bp neuron,draw=NNOutput,fill=NNOutput!18},
  bp link/.style={draw=NNLine,line width=.85pt},
  bp forward/.style={-{Stealth[length=1.7mm,width=1.2mm]},draw=NNHidden,
    line width=.85pt},
  bp backward/.style={-{Stealth[length=1.8mm,width=1.3mm]},draw=NNGradient,
    line width=1.2pt,dashed},
  bp adjoint/.style={text=NNGradient,inner sep=2pt}]
  \foreach \x/\title/\symbol in {
    10/输入/\bm{x},47/预激活/\bm{z}^{(1)},84/隐藏表示/\bm{h}^{(1)},
    121/输出得分/\bm{a},156/损失/\ell}{
    \node[font=\figurefont\bfseries] at (\x,29) {\title};
    \node at (\x,21) {$\symbol$};
  }
  % 先画连接，再画节点，使细线在神经元边缘干净地结束。
  \foreach \iy in {9,0,-9}
    \foreach \hy in {13.5,4.5,-4.5,-13.5}
      \draw[bp link] (10,\iy)--(47,\hy);
  \foreach \hy in {13.5,4.5,-4.5,-13.5}
    \foreach \oy in {6,-6}
      \draw[bp link] (84,\hy)--(121,\oy);
  \foreach \i/\y in {1/9,2/0,3/-9}
    \node[bp input] (x\i) at (10,\y) {$x_{\i}$};
  \foreach \i/\y in {1/13.5,2/4.5,3/-4.5,4/-13.5}{
    \node[bp hidden,fill=white] (z\i) at (47,\y) {$z_{\i}$};
    \node[bp hidden] (h\i) at (84,\y) {$h_{\i}$};
    \draw[bp forward] (z\i)--(h\i);
  }
  \foreach \i/\y in {1/6,2/-6}
    \node[bp output] (a\i) at (121,\y) {$a_{\i}$};
  \node[circle,draw=NNGradient,fill=NNGradient!18,line width=1.1pt,
    minimum size=10mm,inner sep=0pt] (loss) at (156,0) {$\ell$};
  \draw[bp forward,draw=NNOutput] (a1)--(loss.north west);
  \draw[bp forward,draw=NNOutput] (a2)--(loss.south west);
  \node[text=NNInput] at (28.5,15) {$\bm{W}_1,\bm{b}_1$};
  \node[text=NNHidden] at (65.5,15) {$\sigma_1$};
  \node[text=NNHidden] at (102.5,15) {$\bm{W}_2,\bm{b}_2$};
  \node[text=NNOutput] at (140,15) {$\ell(\bm{a},y)$};
  \draw[NNLine,line width=.45pt] (1,-20)--(164,-20);
  \node[bp adjoint] (b1) at (10,-35) {$\nabla_{\bm{x}}\ell$};
  \node[bp adjoint] (b2) at (47,-35) {$\bm{\delta}^{(1)}$};
  \node[bp adjoint] (b3) at (84,-35) {$\nabla_{\bm{h}^{(1)}}\ell$};
  \node[bp adjoint] (b4) at (121,-35) {$\bm{\delta}^{(2)}$};
  \node[bp adjoint] (b5) at (156,-35) {$1$};
  \draw[bp backward] (b2)--node[above] {$\bm{W}_1^{\top}$} (b1);
  \draw[bp backward] (b3)--node[above] {$\odot\sigma_1'$} (b2);
  \draw[bp backward] (b4)--node[above] {$\bm{W}_2^{\top}$} (b3);
  \draw[bp backward] (b5)--node[above] {$\nabla_{\bm{a}}\ell$} (b4);
  \draw[bp backward] (b2.south)--(47,-47);
  \draw[bp backward] (b4.south)--(121,-47);
  \node[align=center,inner sep=0pt] at (47,-55)
    {$\nabla_{\bm{W}_1}\ell=\bm{\delta}^{(1)}\bm{x}^{\top}$\\[2pt]
     $\nabla_{\bm{b}_1}\ell=\bm{\delta}^{(1)}$};
  \node[align=center,inner sep=0pt] at (121,-55)
    {$\nabla_{\bm{W}_2}\ell=\bm{\delta}^{(2)}(\bm{h}^{(1)})^{\top}$\\[2pt]
     $\nabla_{\bm{b}_2}\ell=\bm{\delta}^{(2)}$};
  \draw[bp forward,draw=NNInput] (28,-70)--(38,-70);
  \node[anchor=west] at (41,-70) {前向计算};
  \draw[bp backward] (99,-70)--(89,-70);
  \node[anchor=west] at (102,-70) {梯度回传};
\end{tikzpicture}
